\chapter{Mathematical induction}
\chaptercontents{A & Peano's axioms and recursion (\S4.1) \\ B & Weak induction (\S4.2) \\ C & Strong induction and well-ordering (\S4.3) \\ D & Chapter 4 exercises (4.E)}%
{\S4.1 is not in the planning (covered briefly).\\ \fE\ 4.2: 5, 11\quad \fR\ 4.2: 8, 17, 20\\ \fE\ 4.3: 8, 9, 14, 15\quad \fR\ 4.3: 10, 18\\ \fE\ 4.E: 10, 11, 13, 16--20, 25--27, 29--32\\ \fR\ 4.E: 12, 14, 15, 21, 24, 28, 33--35\\ Tested in the \textbf{exam} (4 Nov).}

\hsection{Peano's axioms and recursion}

\begin{key}[{Definition 4.1.1 and Theorem 4.1.2 \textperiodcentered\ what $\N$ is}]
A \term{notion of natural numbers} is a set $\N$ with a zero $z\in\N$ and a successor function $s:\N\to\N$ such that (i) $z\notin s[\N]$; (ii) $s$ is injective; (iii) $\N$ is generated by $z$ and $s$: for every set $X$, if $z\in X$ and ($n\in X\Rightarrow s(n)\in X$ for all $n\in\N$), then $\N\subseteq X$ (\term{Peano's axioms}). Write $0=z$, $1=s(0)$, $2=s(1)$, \dots. \textbf{Recursion theorem:} for any set $X$, $a\in X$ and $h:\N\times X\to X$ there is a unique $f:\N\to X$ with $f(0)=a$ and $f(s(n))=h(n,f(n))$. \textbf{Strategy 4.1.3 (definition by recursion):} define $f(0)$, then define $f(s(n))$ in terms of $n$ and $f(n)$.
\end{key}
\begin{key}[{Examples 4.1.4, 4.1.7 and Definitions 4.1.10--4.1.15 \textperiodcentered\ the recursive definitions}]
\[
m+0=m,\ \ m+s(n)=s(m+n);\qquad m\cdot0=0,\ \ m\cdot(n+1)=(m\cdot n)+m;\qquad s(n)=n+1\ \text{(Prop.\ 4.1.6)}.
\]
\[
\sum_{k=1}^{0}a_k=0,\ \ \sum_{k=1}^{n+1}a_k=\Big(\sum_{k=1}^na_k\Big)+a_{n+1};\qquad \prod_{k=1}^{0}a_k=1,\ \ \prod_{k=1}^{n+1}a_k=\Big(\prod_{k=1}^na_k\Big)\cdot a_{n+1}.
\]
\[
0!=1,\quad (n+1)!=(n+1)\cdot n!;\qquad \binom n0=1,\quad \binom0{k+1}=0,\quad \binom{n+1}{k+1}=\binom nk+\binom n{k+1}.
\]
(The last three rules generate Pascal's triangle: each entry is the sum of the two above it.)
Model (Proposition 4.1.5, $2+2=4$): $2+2=2+s(1)=s(2+1)=s(2+s(0))=s(s(2+0))=s(s(2))=s(3)=4$, one definition per step.
\end{key}

\begin{bookex}{4.1.13}
Let $x_1,x_2\in\R$. Working strictly from the definitions of indexed sum and product, prove that $\sum_{i=1}^2x_i=x_1+x_2$ and $\prod_{i=1}^2x_i=x_1\cdot x_2$.
\solution
By Definition 4.1.10 (three times): $\sum_{i=1}^2x_i=\big(\sum_{i=1}^1x_i\big)+x_2=\big(\big(\sum_{i=1}^0x_i\big)+x_1\big)+x_2=(0+x_1)+x_2=x_1+x_2$. By Definition 4.1.11: $\prod_{i=1}^2x_i=\big(\prod_{i=1}^1x_i\big)\cdot x_2=\big(\big(\prod_{i=1}^0x_i\big)\cdot x_1\big)\cdot x_2=(1\cdot x_1)\cdot x_2=x_1\cdot x_2$.\qed
\end{bookex}

\begin{trybox}[{4.1.9, 4.1.16 \fO}]
\textbf{4.1.9} Define $m^n$ by recursion on $n$ and prove $2^2=4$ from the recursive definitions.\quad \textbf{4.1.16} Write down the next two rows of Pascal's triangle (after $1\ 5\ 10\ 10\ 5\ 1$).
\end{trybox}

\hsection{Weak induction}

\begin{result}[{Theorem 4.2.1 and Strategy 4.2.2 \textperiodcentered\ proof by (weak) induction}]
To prove $\forall n\ge n_0,\ p(n)$ it suffices to prove: \textbf{(base case)} $p(n_0)$; \textbf{(induction step)} for all $n\ge n_0$, if $p(n)$ then $p(n+1)$. In the step, $p(n)$ is the \term{induction hypothesis} and $p(n+1)$ the \term{induction goal}. (Proof: $X=\{n\in\N\mid p(n_0+n)\}$ contains $0$ and is closed under $n\mapsto n+1$, so $\N\subseteq X$ by Peano axiom (iii).)
\end{result}
\begin{strategy}[{How the book writes an induction proof (page 136)}]
\begin{enumerate}[leftmargin=1.5em,itemsep=1pt]
\item Say ``We proceed by induction on $n\ge n_0$'' and \textbf{state $p(n)$ explicitly}.
\item \textbf{(Base case)} ``We need to prove $p(n_0)$: \dots''.
\item \textbf{(Induction step)} ``Let $n\ge n_0$ and suppose $p(n)$. We need to prove $p(n+1)$: \dots''. Write out the goal before proving it; mark the line where the induction hypothesis is used.
\item ``The result follows by induction.''
\end{enumerate}
Models: Proposition 4.2.3 ($\sum_{k=1}^nk=\frac{n(n+1)}2$: $\sum_{k=1}^{n+1}k=\frac{n(n+1)}2+(n+1)=\frac{(n+1)(n+2)}2$); Proposition 4.2.4 ($3\mid n^3-n$: $(n+1)^3-(n+1)=(n^3-n)+3n^2+3n=3(k+n^2+n)$); Proposition 4.2.6 ($3n<2^n$ for $n\ge4$: $3(n+1)=3n+3<2^n+3<2^n+2^n=2^{n+1}$ since $3<2^4\le2^n$). The proof for $n\ge4$ says nothing about $n<4$: check those separately ($p(0)$ true, $p(1),p(2),p(3)$ false).
\end{strategy}

\begin{bookex}[essential]{4.2.5}
Prove by induction that $\sum_{k=0}^n2^k=2^{n+1}-1$ for all $n\in\N$.
\solution
We proceed by induction on $n\ge0$, with $p(n)$: $\sum_{k=0}^n2^k=2^{n+1}-1$.\\
\emph{(Base case)} $\sum_{k=0}^02^k=2^0=1=2^1-1$. So $p(0)$ holds.\\
\emph{(Induction step)} Let $n\ge0$ and suppose $\sum_{k=0}^n2^k=2^{n+1}-1$. We need to prove $\sum_{k=0}^{n+1}2^k=2^{n+2}-1$. Indeed
\begin{align*}
\sum_{k=0}^{n+1}2^k&=\Big(\sum_{k=0}^n2^k\Big)+2^{n+1}\reason{Definition 4.1.10}\\
&=(2^{n+1}-1)+2^{n+1}\reason{induction hypothesis}\\
&=2\cdot2^{n+1}-1=2^{n+2}-1 .
\end{align*}
The result follows by induction.\qed
\end{bookex}

\begin{bookex}[recommended]{4.2.8}
Write down the statement $p(n)$ that Example 4.2.7 (``all horses have the same colour'') tried to prove for $n\ge1$; write down, with quantifiers, exactly what the induction step must prove; explain why the argument fails.
\solution
$p(n)$: ``every set of $n$ horses has the same colour'' (all horses in the set share one colour). Base case $p(1)$: a set with one horse trivially has one colour: correct. The induction step must prove $\forall n\ge1,\ (p(n)\Rightarrow p(n+1))$: \emph{for every} $n\ge1$, if every set of $n$ horses is monochrome then every set of $n+1$ horses is. The argument removes the first horse and the last horse from a set $X$ of $n+1$ horses and uses that the two resulting $n$-sets overlap, so that the colour of the overlap forces all $n+1$ to agree. For $n=1$ the two $n$-sets $\{h_2\}$ and $\{h_1\}$ do \emph{not} overlap, so nothing links the colour of $h_1$ to that of $h_2$: the implication $p(1)\Rightarrow p(2)$ was never proved. The step works for $n\ge2$ only, and the chain from the base case is broken at its first link.\qed
\end{bookex}

\begin{bookex}[essential]{4.2.11}
Let $a,r\in\R$ with $r\ne1$. Prove that $\sum_{k=0}^nar^k=\dfrac{a(1-r^{n+1})}{1-r}$ for all $n\in\N$.
\solution
We proceed by induction on $n\ge0$.\\
\emph{(Base case)} $\sum_{k=0}^0ar^k=ar^0=a$, and $\frac{a(1-r^1)}{1-r}=a$ since $r\ne1$. So the formula holds for $n=0$.\\
\emph{(Induction step)} Let $n\ge0$ and suppose $\sum_{k=0}^nar^k=\frac{a(1-r^{n+1})}{1-r}$. We need to prove $\sum_{k=0}^{n+1}ar^k=\frac{a(1-r^{n+2})}{1-r}$. Now
\begin{align*}
\sum_{k=0}^{n+1}ar^k&=\sum_{k=0}^nar^k+ar^{n+1}\reason{Definition 4.1.10}\\
&=\frac{a(1-r^{n+1})}{1-r}+ar^{n+1}\reason{induction hypothesis}\\
&=\frac{a(1-r^{n+1})+ar^{n+1}(1-r)}{1-r}=\frac{a-ar^{n+1}+ar^{n+1}-ar^{n+2}}{1-r}=\frac{a(1-r^{n+2})}{1-r}.
\end{align*} The result follows by induction. (Exercise 4.2.5 is the case $a=1$, $r=2$.)\qed
\end{bookex}

\begin{trybox}[{4.2.12 \fO}]
\textbf{4.2.12} For sets $X_k$ ($k\in[n+1]$), prove by induction on $n$ that there is a bijection $\prod_{k\in[n+1]}X_k\to\big(\prod_{k\in[n]}X_k\big)\times X_n$. \hint{elements of the left are $(a_0,\dots,a_n)$; of the right, $((a_0,\dots,a_{n-1}),a_n)$.}
\end{trybox}

\hsub{Binomials and the binomial theorem}
\begin{result}[{Theorem 4.2.15, Example 4.2.16, Theorem 4.2.18}]
$\binom nk=\dfrac{n!}{k!(n-k)!}$ if $k\le n$, and $\binom nk=0$ if $k>n$ (proved from the recursion by induction on $n$, with four cases in the step). Consequences: $\binom n0=\binom nn=1$; $\binom nk\binom k\ell=\binom n\ell\binom{n-\ell}{k-\ell}$ for $\ell\le k\le n$. \textbf{Binomial theorem:} $(x+y)^n=\sum_{k=0}^n\binom nkx^ky^{n-k}$ for $n\in\N$, $x,y\in\R$ (proved for $(1+x)^n$ by induction, reindexing and using the Pascal recursion, then scaling by $y^n$). Example 4.2.13/4.2.19: $\sum_{k=0}^n\binom nk=2^n$, by induction or at once from $(1+1)^n$.
\end{result}

\begin{bookex}[recommended]{4.2.17}
Prove that $\binom nk=\binom n{n-k}$ for all $n,k\in\N$ with $k\le n$.
\solution
Since $k\le n$ we have $n-k\in\N$ and $n-k\le n$, so Theorem 4.2.15 applies to both sides:
\[
\binom n{n-k}=\frac{n!}{(n-k)!\,(n-(n-k))!}=\frac{n!}{(n-k)!\,k!}=\frac{n!}{k!\,(n-k)!}=\binom nk .
\]\qed
\end{bookex}

\begin{bookex}[recommended]{4.2.20}
Give a proof of $\sum_{i=0}^n(-1)^i\binom ni=0$ using the binomial theorem.
\solution
For $n\ge1$, apply Theorem 4.2.18 with $x=-1$, $y=1$: $\sum_{i=0}^n\binom ni(-1)^i1^{n-i}=(-1+1)^n=0^n=0$ (the hint: $(-1)^i\binom ni=\binom ni(-1)^i1^{n-i}$). Note that the identity needs $n\ge1$: for $n=0$ the sum is $(-1)^0\binom00=1$, matching $0^0=1$.\qed
\end{bookex}

\begin{trybox}[{4.2.14 \fO}]
\textbf{4.2.14} Prove $\sum_{i=0}^n(-1)^i\binom ni=0$ (for $n\ge1$) by induction, in the style of Example 4.2.13.
\end{trybox}

\hsection{Strong induction and well-ordering}

\begin{result}[{Theorems 4.3.2, 4.3.5 and Strategies 4.3.3, 4.3.6 \textperiodcentered\ strong induction}]
To prove $\forall n\ge n_0,\ p(n)$ it suffices to prove $p(n_0)$ and: for all $n\ge n_0$, if $p(k)$ holds \emph{for all $n_0\le k\le n$}, then $p(n+1)$. \textbf{Multiple base cases:} prove $p(n_0),\dots,p(n_1)$ and: for all $n\ge n_1$, if $p(k)$ for all $n_0\le k\le n$ then $p(n+1)$. Strong induction is no stronger than weak induction (it is proved from it, applied to $q(n)$: ``$p(k)$ for all $n_0\le k\le n$''); use it when the goal $p(n+1)$ depends on earlier terms than $p(n)$ (Example 4.3.1: $b_{n+1}=1+\sum_{k\le n}b_k$ needs all $b_k$). \textbf{Get the quantifiers right:} ``Fix $n\ge n_1$ and assume $p(k)$ for all $n_0\le k\le n$.''
\end{result}
\begin{strategy}[{Model: Example 4.3.7}]
$a_0=0$, $a_1=1$, $a_n=3a_{n-1}-2a_{n-2}$ ($n\ge2$); claim $a_n=2^n-1$. \emph{Base cases:} $a_0=0=2^0-1$, $a_1=1=2^1-1$. \emph{Step:} fix $n\ge1$ and assume $a_k=2^k-1$ for all $0\le k\le n$. Since $n+1\ge2$ the recursion applies: $a_{n+1}=3a_n-2a_{n-1}=3(2^n-1)-2(2^{n-1}-1)=3\cdot2^n-3-2^n+2=2\cdot2^n-1=2^{n+1}-1$. The step uses $a_n$ \emph{and} $a_{n-1}$, which is why there are two base cases and why $n\ge1$ (not $n\ge0$) is fixed in the step.
\end{strategy}

\begin{bookex}[essential]{4.3.8}
Define $a_0=4$, $a_1=9$ and $a_n=5a_{n-1}-6a_{n-2}$ for $n\ge2$. Prove that $a_n=3\cdot2^n+3^n$ for all $n\in\N$.
\solution
We proceed by strong induction on $n\ge0$ with two base cases.\\
\emph{(Base cases)} $a_0=4=3\cdot2^0+3^0$ and $a_1=9=3\cdot2^1+3^1=6+3$.\\
\emph{(Induction step)} Fix $n\ge1$ and assume $a_k=3\cdot2^k+3^k$ for all $0\le k\le n$. We need to prove $a_{n+1}=3\cdot2^{n+1}+3^{n+1}$. Since $n+1\ge2$,
\begin{align*}
a_{n+1}&=5a_n-6a_{n-1}\reason{definition}\\
&=5(3\cdot2^n+3^n)-6(3\cdot2^{n-1}+3^{n-1})\reason{induction hypothesis for $k=n,n-1$}\\
&=15\cdot2^n+5\cdot3^n-9\cdot2^n-2\cdot3^n\reason{$6\cdot2^{n-1}=3\cdot2^n$, $6\cdot3^{n-1}=2\cdot3^n$}\\
&=6\cdot2^n+3\cdot3^n=3\cdot2^{n+1}+3^{n+1}.
\end{align*}
The result follows by induction.\qed
\end{bookex}

\begin{bookex}[essential]{4.3.9}
The Tribonacci sequence is $t_0=0$, $t_1=0$, $t_2=1$, $t_n=t_{n-1}+t_{n-2}+t_{n-3}$ for $n\ge3$. Prove that $t_n\le2^{n-3}$ for all $n\ge3$.
\solution
We proceed by strong induction on $n\ge3$ with base cases $n=3,4,5$ (the step uses three previous terms).\\
\emph{(Base cases)} $t_3=1+0+0=1\le2^0$; $t_4=1+1+0=2\le2^1$; $t_5=2+1+1=4\le2^2$.\\
\emph{(Induction step)} Fix $n\ge5$ and assume $t_k\le2^{k-3}$ for all $3\le k\le n$. We need $t_{n+1}\le2^{n-2}$. Since $n+1\ge6$, $t_{n+1}=t_n+t_{n-1}+t_{n-2}$, and $n,n-1,n-2\ge3$, so the hypothesis applies to all three:
\[
t_{n+1}\le2^{n-3}+2^{n-4}+2^{n-5}=2^{n-5}(4+2+1)=7\cdot2^{n-5}<8\cdot2^{n-5}=2^{n-2}.
\]
The result follows by induction.\qed
\end{bookex}

\begin{bookex}[recommended]{4.3.10}
Prove that for all $n\ge8$, a value of $n$ dubloons can be obtained using only $3$- and $5$-dubloon coins.
\solution
Let $p(n)$: ``$n=3a+5b$ for some $a,b\in\N$''. We proceed by strong induction on $n\ge8$ with base cases $8,9,10$ (hint: if $p(n)$ then $p(n+3)$, by adding one $3$-coin).\\
\emph{(Base cases)} $8=3+5$; $9=3+3+3$; $10=5+5$.\\
\emph{(Induction step)} Fix $n\ge10$ and assume $p(k)$ for all $8\le k\le n$. We need $p(n+1)$. Since $n\ge10$, $n-2\ge8$, so $p(n-2)$ holds: $n-2=3a+5b$. Then $n+1=3(a+1)+5b$, so $p(n+1)$ holds.\\
The result follows by induction. (For $n=7$: $7=3a+5b$ would need $b\in\{0,1\}$, giving $3a=7$ or $3a=2$, impossible.)\qed
\end{bookex}

\hsub{The well-ordering principle}
\begin{result}[{Theorem 4.3.11, Lemma 4.3.13, Proposition 4.3.12}]
\textbf{Well-ordering principle:} every non-empty set of natural numbers has a least element (proved by strong induction on an element $n$ of the set). \textbf{Lemma 4.3.13:} every positive rational $q$ is $\frac ab$ with $a,b\in\N$ non-zero and $1$ the only natural number dividing both (take the least $a$ with $q=\frac ab$; a common divisor $d\ge2$ gives $q=\frac{a'}{b'}$ with $a'=\frac ad<a$). \textbf{$\sqrt2$ is irrational:} if $\sqrt2=\frac ab$ in lowest terms then $a^2=2b^2$, so $a$ is even by Proposition 1.1.63 (``$a^2$ even $\Rightarrow$ $a$ even''), $a=2c$, $2c^2=b^2$, so $b$ is even too; then $2$ divides both $a$ and $b$, contradiction. \emph{Writing tip:} prove a repeated step once as a lemma and quote it.
\end{result}

\begin{bookex}[essential]{4.3.14}
What goes wrong in the proof of Proposition 4.3.12 if we try to prove that $\sqrt4$ is irrational?
\solution
Suppose $\sqrt4=\frac ab$ in lowest terms. Squaring gives $a^2=4b^2$, so $a^2$ is even and hence $a$ is even, $a=2c$. Substituting: $4c^2=4b^2$, so $c^2=b^2$. At this point the argument for $\sqrt2$ deduced ``$b^2$ is even''; here we only get $c^2=b^2$, which says nothing about the parity of $b$. No contradiction arises, and indeed $\sqrt4=\frac21$ with $a=2$, $b=1$ satisfies everything. The step that fails is the one that depended on the prime $2$ dividing the right-hand side an odd number of times.\qed
\end{bookex}

\begin{bookex}[essential]{4.3.15}
Prove that $\sqrt3$ is irrational.
\solution
\emph{Lemma.} If $a\in\Z$ and $3$ divides $a^2$, then $3$ divides $a$. \emph{Proof.} By the division theorem $a=3q+r$ with $r\in\{0,1,2\}$. If $r=1$ then $a^2=9q^2+6q+1=3(3q^2+2q)+1$; if $r=2$ then $a^2=9q^2+12q+4=3(3q^2+4q+1)+1$. In both cases $a^2$ leaves remainder $1$, so $3\nmid a^2$ (remainders are unique, Theorem 0.42). Hence $3\mid a^2$ forces $r=0$, i.e.\ $3\mid a$.\\
\emph{Proof of the claim.} Suppose $\sqrt3$ is rational. Since $\sqrt3>0$, by Lemma 4.3.13 we can write $\sqrt3=\frac ab$ with $a,b$ positive natural numbers whose only common natural divisor is $1$. Squaring: $a^2=3b^2$, so $3\mid a^2$ and by the Lemma $3\mid a$: $a=3c$. Then $9c^2=3b^2$, so $b^2=3c^2$, $3\mid b^2$, and $3\mid b$ by the Lemma. So $3$ divides both $a$ and $b$, contradicting lowest terms. Hence $\sqrt3$ is irrational.\qed
\end{bookex}

\hsub{Infinite descent}
\begin{key}[{Definition 4.3.16 and Strategy 4.3.19}]
$X\subseteq\R$ has \term{infinite descent} if for all $x\in X$ there is $y\in X$ with $y<x$ (Example 4.3.17: $\Z$). The well-ordering principle says no \emph{non-empty subset of $\N$} has infinite descent. \textbf{Proof by infinite descent:} to derive a contradiction from $p$, use $p$ to build a non-empty $X\subseteq\N$ with infinite descent. Typical use: ``$X$ = the set of counterexamples'' (Example 4.3.21, existence of prime factorisation: a smallest counterexample $n=ab$ yields a smaller one) or ``$X$ = the set of candidate remainders'' (Example 4.3.20: if the least $r$ with $a=qb+r$ had $r\ge b$, then $r-b$ would be a smaller candidate).
\end{key}

\begin{bookex}[recommended]{4.3.18}
Prove that the interval $(0,1)$ has infinite descent.
\solution
Let $x\in(0,1)$. Put $y=\frac x2$. Then $0<y$ (as $x>0$), $y<x$ (as $x>0$), and $y<x<1$, so $y\in(0,1)$ and $y<x$. Hence $(0,1)$ has infinite descent. (No contradiction follows: $(0,1)\nsubseteq\N$.)\qed
\end{bookex}

\begin{trybox}[{4.3.22 \fO}]
\textbf{4.3.22} Prove by infinite descent: if $\sqrt n$ ($n\in\N$) is not an integer then it is irrational. \hint{if $\sqrt n=\frac ab$, then $a(\sqrt n-\lfloor\sqrt n\rfloor)$ and $b(\sqrt n-\lfloor\sqrt n\rfloor)$ are integers; consider $\sqrt n\cdot\frac{\sqrt n-\lfloor\sqrt n\rfloor}{\sqrt n-\lfloor\sqrt n\rfloor}$.}
\end{trybox}

\begin{warn}
\begin{itemize}
\item The induction step must work for \emph{every} $n\ge n_0$, including $n=n_0$ (the horses, 4.2.8). When the step uses $p(n-1)$ or $p(n-2)$ you need extra base cases and must fix $n\ge n_1$.
\item Write the induction goal before proving it and never assume it; mark where the induction hypothesis is used.
\item Divisibility statements (``$3\mid n^3-n$'', 4.E.16, 4.E.17): unfold ``$=3k$ for some $k\in\Z$'' and exhibit the new quotient explicitly.
\item ``$3n<2^n$ for $n\ge4$'' says nothing about $n<4$: check small cases one by one (4.E.10).
\end{itemize}
\end{warn}

\hsection{Chapter 4 exercises}

\hsub{Recursive definitions (4.E.1--4.E.9)}
\begin{trybox}[{4.E.1--4.E.9 \fO}]
\textbf{4.E.1--4.E.5} From the recursive definitions: $1+3=4$; $0+5=5$; $2\cdot3=6$; $0\cdot5=0$; $2^3=8$.\quad \textbf{4.E.6} Define $\exists^{=n}$ by recursion on $n$.\quad \textbf{4.E.7} From Definition 4.1.15, $\binom42=6$.\quad \textbf{4.E.8} Trailing zeros of $41!$ in decimal and in binary. \hint{the number of trailing zeros in base $b$ is the largest $r$ with $b^r\mid n$.}\quad \textbf{4.E.9} $(N,z,s)$ is a notion of natural numbers iff for every $X$, $a\in X$, $f:X\to X$ there is a unique $h:N\to X$ with $h(z)=a$ and $f\circ h=h\circ s$.
\end{trybox}

\hsub{Proofs by induction (4.E.10--4.E.17)}
\begin{bookex}[essential]{4.E.10}
Find all natural numbers $n$ such that $n^5<5^n$.
\solution
Small cases: $0<1$ (true); $1<5$ (true); $32<25$ (false); $243<125$ (false); $1024<625$ (false); $3125<3125$ (false); $7776<15625$ (true). \textbf{Claim:} $n^5<5^n$ for all $n\ge6$. By induction on $n\ge6$. \emph{Base case:} $6^5=7776<15625=5^6$. \emph{Step:} let $n\ge6$ and suppose $n^5<5^n$. Then $\big(\frac{n+1}n\big)^5=\big(1+\frac1n\big)^5\le\big(\frac76\big)^5=\frac{16807}{7776}<5$, so $(n+1)^5<5n^5<5\cdot5^n=5^{n+1}$. The result follows by induction.\\
Answer: $n^5<5^n$ exactly for $n\in\{0,1\}\cup\{n\in\N\mid n\ge6\}$.\qed
\end{bookex}

\begin{bookex}[essential]{4.E.11}
Prove that $(1+x)^{123\,456\,789}\ge1+123\,456\,789\,x$ for all real $x\ge-1$.
\solution
Nothing is special about $123\,456\,789$ (the hint): we prove $(1+x)^n\ge1+nx$ for all $n\in\N$ and all real $x\ge-1$, by induction on $n$ (not on $x$: induction only reaches natural numbers). Fix $x\ge-1$. \emph{Base case:} $(1+x)^0=1\ge1+0\cdot x$. \emph{Step:} let $n\ge0$ and suppose $(1+x)^n\ge1+nx$. Since $1+x\ge0$ we may multiply the hypothesis by $1+x$ without changing the inequality:
\[
(1+x)^{n+1}=(1+x)(1+x)^n\ge(1+x)(1+nx)=1+(n+1)x+nx^2\ge1+(n+1)x,
\]
as $nx^2\ge0$. By induction the inequality holds for all $n$; take $n=123\,456\,789$.\qed
\end{bookex}

\begin{bookex}[recommended]{4.E.12}
Let $a\in\N$ end in the digit $6$. Prove that $a^n$ ends in $6$ for all $n\ge1$.
\solution
``Ends in $6$'' means $a=10q+6$ for some $q\in\N$ (the last decimal digit is the remainder on division by $10$). Induction on $n\ge1$. \emph{Base case:} $a^1=a$ ends in $6$. \emph{Step:} let $n\ge1$ and suppose $a^n=10m+6$. Then $a^{n+1}=a^n\cdot a=(10m+6)(10q+6)=100mq+60m+60q+36=10(10mq+6m+6q+3)+6$, which ends in $6$. The result follows by induction.\qed
\end{bookex}

\begin{bookex}[essential]{4.E.13}
Let $a,b\in\R$ and let $a_0=a$, $a_1=b$ and $a_n=\dfrac{a_{n-1}+a_{n+1}}2$ for all $n\ge1$. Prove that $a_n=a+(b-a)n$ for all $n\in\N$.
\solution
Rewrite the recursion as $a_{n+1}=2a_n-a_{n-1}$ for $n\ge1$. We proceed by strong induction with base cases $n=0,1$. \emph{Base cases:} $a_0=a=a+(b-a)\cdot0$ and $a_1=b=a+(b-a)\cdot1$. \emph{Step:} fix $n\ge1$ and assume $a_k=a+(b-a)k$ for all $0\le k\le n$. Then
\[
a_{n+1}=2a_n-a_{n-1}=2\big(a+(b-a)n\big)-\big(a+(b-a)(n-1)\big)=a+(b-a)(2n-n+1)=a+(b-a)(n+1).
\]
The result follows by induction.\qed
\end{bookex}

\begin{bookex}[recommended]{4.E.14}
Let $f:\R\to\R$ satisfy $f(x+y)=f(x)+f(y)$ for all $x,y$. (a) Prove by induction that there is $a\in\R$ with $f(n)=an$ for all $n\in\N$. (b) Deduce $f(n)=an$ for all $n\in\Z$. (c) Deduce $f(x)=ax$ for all $x\in\Q$.
\solution
\emph{(a)} Put $a=f(1)$. First $f(0)=f(0+0)=2f(0)$, so $f(0)=0=a\cdot0$: base case. Step: if $f(n)=an$ then $f(n+1)=f(n)+f(1)=an+a=a(n+1)$. By induction $f(n)=an$ for all $n\in\N$.\\
\emph{(b)} For $n\in\N$: $0=f(0)=f(n+(-n))=f(n)+f(-n)=an+f(-n)$, so $f(-n)=-an=a(-n)$. Every integer is $n$ or $-n$ for some $n\in\N$.\\
\emph{(c)} Let $x\in\Q$ and write $x=\frac mb$ with $m\in\Z$ and $b\ge1$. First, $f(bx)=bf(x)$ for all $b\in\N$, by induction on $b$: $f(0\cdot x)=f(0)=0$, and $f((b+1)x)=f(bx+x)=f(bx)+f(x)=bf(x)+f(x)$. Hence $bf(x)=f(bx)=f(m)=am$ by (b), so $f(x)=\frac{am}b=ax$.\qed
\end{bookex}

\begin{bookex}[recommended]{4.E.15}
Let $f:\R\to\R$ with $f(0)>0$ and $f(x+y)=f(x)f(y)$ for all $x,y$. Prove there is a positive real $a$ with $f(x)=a^x$ for all $x\in\Q$.
\solution
$f(0)=f(0)f(0)$ with $f(0)>0$ gives $f(0)=1$. For any $x$, $f(x)f(-x)=f(0)=1$, so $f(x)\ne0$, and $f(x)=f(\frac x2)^2>0$. Put $a=f(1)>0$. \emph{Naturals:} $f(0)=1=a^0$; if $f(n)=a^n$ then $f(n+1)=f(n)f(1)=a^{n+1}$; by induction $f(n)=a^n$ for $n\in\N$. \emph{Integers:} $f(-n)=1/f(n)=a^{-n}$. \emph{Rationals:} for $x=\frac mb$ with $m\in\Z$, $b\ge1$: by induction on $b$, $f(bx)=f(x)^b$; so $f(x)^b=f(m)=a^m$, and since $f(x)>0$, $f(x)$ is the unique positive $b$-th root of $a^m$, namely $a^{m/b}=a^x$.\qed
\end{bookex}

\begin{bookex}[essential]{4.E.16}
Let $a,b\in\Z$. Prove that $b-a$ divides $b^n-a^n$ for all $n\in\N$.
\solution
Induction on $n\ge0$; $p(n)$: $b^n-a^n=(b-a)k$ for some $k\in\Z$. \emph{Base case:} $b^0-a^0=0=(b-a)\cdot0$. \emph{Step:} let $n\ge0$ and suppose $b^n-a^n=(b-a)k$. Then
\[
b^{n+1}-a^{n+1}=b\cdot b^n-a\cdot a^n=b(b^n-a^n)+(b-a)a^n=(b-a)(bk+a^n),
\]
and $bk+a^n\in\Z$. So $b-a\mid b^{n+1}-a^{n+1}$. The result follows by induction.\qed
\end{bookex}

\begin{bookex}[essential]{4.E.17}
Prove by induction that $7^n-2\cdot4^n+1$ is divisible by $18$ for all $n\in\N$.
\solution
The book's hint: $7^{n+1}-2\cdot4^{n+1}+1=4(7^n-2\cdot4^n+1)+3(7^n-1)$ for all $n$ (check: the right side is $7\cdot7^n-8\cdot4^n+4-3=7^{n+1}-2\cdot4^{n+1}+1$). So we need a separate fact about $7^n-1$.\\
\emph{Lemma: $6\mid7^n-1$ for all $n\in\N$.} By induction: $7^0-1=0=6\cdot0$; and if $7^n-1=6m$ then $7^{n+1}-1=7(7^n-1)+6=6(7m+1)$.\\
\emph{Main proof}, by induction on $n\ge0$. \emph{Base case:} $7^0-2\cdot4^0+1=1-2+1=0=18\cdot0$. \emph{Step:} let $n\ge0$ and suppose $7^n-2\cdot4^n+1=18k$ with $k\in\Z$. By the Lemma, $7^n-1=6m$ for some $m\in\Z$. Then
\[
7^{n+1}-2\cdot4^{n+1}+1=4\cdot18k+3\cdot6m=18(4k+m),
\]
so $18$ divides $7^{n+1}-2\cdot4^{n+1}+1$. The result follows by induction.\qed
\end{bookex}

\hsub{Some examples from calculus (4.E.18--4.E.23)}
\begin{bookex}[essential]{4.E.18}
Prove that $\dfrac{d^n}{dx^n}(xe^x)=(n+x)e^x$ for all $n\in\N$.
\solution
Induction on $n$. \emph{Base case:} the $0$-th derivative is the function itself: $xe^x=(0+x)e^x$. \emph{Step:} suppose $\frac{d^n}{dx^n}(xe^x)=(n+x)e^x$. Differentiating once more (product rule): $\frac{d^{n+1}}{dx^{n+1}}(xe^x)=\frac d{dx}\big((n+x)e^x\big)=1\cdot e^x+(n+x)e^x=(n+1+x)e^x$. The result follows by induction.\qed
\end{bookex}

\begin{bookex}[essential]{4.E.19}
Find and prove an expression for $\dfrac{d^n}{dx^n}(x^2e^x)$ for all $n\in\N$.
\solution
Compute: $n=0$: $x^2e^x$; $n=1$: $(x^2+2x)e^x$; $n=2$: $(x^2+4x+2)e^x$; $n=3$: $(x^2+6x+6)e^x$. \textbf{Claim:} $\frac{d^n}{dx^n}(x^2e^x)=\big(x^2+2nx+n(n-1)\big)e^x$. \emph{Base case} $n=0$: $(x^2+0+0)e^x$. \emph{Step:} if the claim holds for $n$, then
\begin{align*}
\frac d{dx}\Big(\big(x^2+2nx+n(n-1)\big)e^x\Big)&=(2x+2n)e^x+\big(x^2+2nx+n(n-1)\big)e^x\\
&=\big(x^2+2(n+1)x+n(n-1)+2n\big)e^x,
\end{align*}
and $n(n-1)+2n=n(n+1)=(n+1)n$, which is the claim for $n+1$. By induction the formula holds for all $n$.\qed
\end{bookex}

\begin{bookex}[essential]{4.E.20}
Let $f,g$ be differentiable (as often as needed). Prove that $\dfrac{d^n}{dx^n}\big(f(x)g(x)\big)=\sum_{k=0}^n\binom nkf^{(k)}(x)g^{(n-k)}(x)$ for all $n\in\N$.
\solution
Induction on $n$. \emph{Base case:} $n=0$: $\binom00f^{(0)}g^{(0)}=fg$. \emph{Step:} suppose the formula for $n$. Differentiate term by term with the product rule:
\[
\frac{d^{n+1}}{dx^{n+1}}(fg)=\sum_{k=0}^n\binom nk\big(f^{(k+1)}g^{(n-k)}+f^{(k)}g^{(n+1-k)}\big)=\sum_{k=1}^{n+1}\binom n{k-1}f^{(k)}g^{(n+1-k)}+\sum_{k=0}^n\binom nkf^{(k)}g^{(n+1-k)}
\]
(reindexing the first sum with $k\to k-1$). Split off $k=n+1$ from the first sum and $k=0$ from the second, and combine the rest with Pascal's rule $\binom n{k-1}+\binom nk=\binom{n+1}k$ (Definition 4.1.15):
\[
=\binom nnf^{(n+1)}g^{(0)}+\sum_{k=1}^n\binom{n+1}kf^{(k)}g^{(n+1-k)}+\binom n0f^{(0)}g^{(n+1)}=\sum_{k=0}^{n+1}\binom{n+1}kf^{(k)}g^{(n+1-k)},
\]
since $\binom nn=1=\binom{n+1}{n+1}$ and $\binom n0=1=\binom{n+1}0$. This is exactly the induction in the proof of the binomial theorem (Theorem 4.2.18), with derivatives in place of powers.\qed
\end{bookex}

\begin{bookex}[recommended]{4.E.21}
Find and prove an expression for the $n$-th derivative of $\log x$ ($x>0$) for all $n\in\N$.
\solution
$n=0$: $\log x$; $n=1$: $x^{-1}$; $n=2$: $-x^{-2}$; $n=3$: $2x^{-3}$; $n=4$: $-6x^{-4}$. \textbf{Claim:} for $n\ge1$, $\frac{d^n}{dx^n}\log x=(-1)^{n-1}(n-1)!\,x^{-n}$. \emph{Base case} $n=1$: $(-1)^00!\,x^{-1}=\frac1x$. \emph{Step:} if the claim holds for $n\ge1$ then $\frac d{dx}\big((-1)^{n-1}(n-1)!\,x^{-n}\big)=(-1)^{n-1}(n-1)!\cdot(-n)x^{-n-1}=(-1)^nn!\,x^{-(n+1)}$, the claim for $n+1$. By induction the formula holds for all $n\ge1$ (and for $n=0$ the derivative is $\log x$ itself).\qed
\end{bookex}

\begin{trybox}[{4.E.22, 4.E.23 \fO}]
\textbf{4.E.22} $(\cos\theta+i\sin\theta)^n=\cos(n\theta)+i\sin(n\theta)$ for all $n\in\N$. \hint{angle addition formulae.}\quad \textbf{4.E.23} (a) $\int_0^{\pi/2}\sin^{n+2}x\,dx=\frac{n+1}{n+2}\int_0^{\pi/2}\sin^nx\,dx$ (integration by parts); (b) $\int_0^{\pi/2}\sin^{2n}x\,dx=\frac{\pi}{2^{2n+1}}\binom{2n}n$ by induction; (c) find $\int_0^{\pi/2}\sin^{2n+1}x\,dx$.
\end{trybox}

\hsub{True--false (4.E.24--4.E.32)}
\begin{bookex}[recommended]{4.E.24}
Every factorial is positive.
\solution
\emph{True.} By induction on $n$: $0!=1>0$; if $n!>0$ then $(n+1)!=(n+1)\cdot n!>0$ as a product of positive numbers. So $n!>0$ for all $n\in\N$.\qed
\end{bookex}

\begin{bookex}[essential]{4.E.25}
Every binomial coefficient is positive.
\solution
\emph{False.} $\binom01=0$ by Definition 4.1.15 (the clause $\binom0{k+1}=0$); more generally $\binom nk=0$ whenever $k>n$ (Theorem 4.2.15).\qed
\end{bookex}

\begin{bookex}[essential]{4.E.26}
Every proof by strong induction can be turned into a proof by weak induction by changing only the induction hypothesis.
\solution
\emph{False.} A proof by strong induction may use $p(k)$ for $k<n$ in its step (Example 4.3.4 uses all of $b_0,\dots,b_n$; Example 4.3.7 uses $a_{n-1}$). Replacing the hypothesis by just $p(n)$ leaves those uses unjustified, so the step no longer proves $p(n+1)$: the rest of the proof must change too (e.g.\ by proving the stronger statement $q(n)$: ``$p(k)$ for all $k\le n$'', as in the proof of Theorem 4.3.2). (Strong and weak induction prove the same statements, but not by the same proof.)\qed
\end{bookex}

\begin{bookex}[essential]{4.E.27}
Every proof by weak induction can be turned into a proof by strong induction by changing only the induction hypothesis.
\solution
\emph{True.} Replace the hypothesis ``$p(n)$'' by ``$p(k)$ for all $n_0\le k\le n$''. The new hypothesis contains the old one (take $k=n$), so every step of the original argument remains valid and still derives $p(n+1)$; the base case is unchanged. By Theorem 4.3.2 the conclusion follows.\qed
\end{bookex}

\begin{bookex}[recommended]{4.E.28}
Let $f:\N\to X$ be a surjection and $p(x)$ a formula on $X$. If $p(f(0))$ and $\forall n\in\N,\ (p(f(n))\Rightarrow p(f(n+1)))$, then $\forall x\in X,\ p(x)$.
\solution
\emph{True.} Let $q(n)$ be $p(f(n))$, a formula with free variable $n\in\N$. The hypotheses say $q(0)$ and $q(n)\Rightarrow q(n+1)$ for all $n$, so $q(n)$ holds for all $n\in\N$ by weak induction. Now let $x\in X$; since $f$ is surjective, $x=f(n)$ for some $n$, and $p(x)=p(f(n))=q(n)$ is true.\qed
\end{bookex}

\begin{bookex}[essential]{4.E.29}
If for all $x\in\N$ there is some $y\le x$ with $p(y)$ false, then $p(x)$ is false for all $x\in\N$.
\solution
\emph{False.} Take $p(x)$: ``$x\ge1$''. The hypothesis holds: for every $x\in\N$, $y=0\le x$ and $p(0)$ is false. But the conclusion fails, since $p(1)$ is true. (The hypothesis only guarantees that $p(0)$ is false: it is not an induction principle.)\qed
\end{bookex}

\begin{bookex}[essential]{4.E.30}
Every non-empty subset of $\Q_{\ge0}$ has a least element.
\solution
\emph{False.} $\{\frac1n\mid n\ge1\}\subseteq\Q_{\ge0}$ is non-empty and has no least element: for any $\frac1n$ in it, $\frac1{n+1}<\frac1n$ is also in it (infinite descent, Definition 4.3.16).\qed
\end{bookex}

\begin{bookex}[essential]{4.E.31}
Every non-empty subset of $\Z$ has a least element.
\solution
\emph{False.} $\Z$ itself is non-empty and has no least element: for $n\in\Z$, $n-1\in\Z$ and $n-1<n$ (Example 4.3.17).\qed
\end{bookex}

\begin{bookex}[essential]{4.E.32}
Every non-empty subset of $\Z^-$ (the negative integers) has a greatest element.
\solution
\emph{True.} Let $X\subseteq\Z^-$ be non-empty and put $Y=\{-x-1\mid x\in X\}$. For $x\le-1$, $-x-1\ge0$, so $Y\subseteq\N$, and $Y$ is non-empty. By the well-ordering principle $Y$ has a least element $-x_0-1$ with $x_0\in X$. For any $x\in X$: $-x-1\ge-x_0-1$, so $x\le x_0$. Hence $x_0$ is the greatest element of $X$.\qed
\end{bookex}

\hsub{Always--sometimes--never (4.E.33--4.E.35)}
\begin{bookex}[recommended]{4.E.33}
Let $n,k,\ell\in\N$ (the book writes $\Z$, but binomial coefficients are defined for natural numbers). Then $\binom n{k+\ell}=\binom nk+\binom n\ell$.
\solution
\emph{Sometimes.} True for $n=0$, $k=\ell=1$: $\binom02=0=\binom01+\binom01$. False for $n=1$, $k=0$, $\ell=1$: $\binom11=1$ but $\binom10+\binom11=2$.\qed
\end{bookex}

\begin{bookex}[recommended]{4.E.34}
Let $m,n\ge2$. Then $(m+n)!=m!+n!$.
\solution
\emph{Never.} First, $j!\le\ell!$ whenever $j\le\ell$ (each step multiplies by a number $\ge1$). Since $m+n-1\ge m$ and $\ge n$, $(m+n-1)!\ge\max(m!,n!)$, so
\[
(m+n)!=(m+n)\cdot(m+n-1)!\ge4\max(m!,n!)>2\max(m!,n!)\ge m!+n!,
\]
using $m+n\ge4$ and $\max(m!,n!)>0$ (4.E.24). So $(m+n)!>m!+n!$ always.\qed
\end{bookex}

\begin{bookex}[recommended]{4.E.35}
Let $p(x)$ be a formula on $\Z$ with $p(0)$ true and, for all $n\in\Z$, $p(n)\Rightarrow p(n+1)$ and $p(n)\Rightarrow p(n-2)$. Then $\forall n\in\Z,\ p(n)$.
\solution
\emph{Always.} By weak induction, $p(n)$ for all $n\in\N$ (base $p(0)$; step $p(n)\Rightarrow p(n+1)$). For negative integers prove $q(m)$: ``$p(-m)$'' for all $m\in\N$ by strong induction with base cases $m=0,1$: $p(0)$ holds, and $p(-1)$ follows from $p(1)$ via $p(1)\Rightarrow p(1-2)$. Step: fix $m\ge1$ and assume $p(-k)$ for all $0\le k\le m$; then $p(-(m-1))$ holds, and $p(-(m-1))\Rightarrow p(-(m-1)-2)=p(-(m+1))$. So $p(n)$ holds for every integer $n$.\qed
\end{bookex}
